\section*{Outline and motivation}
In this final lecture on seismology we consider the free oscillations of the Earth in further detail. Last time we established the 
general theory, and saw how eigenfrequencies can be observed from the Fourier transforms of sufficiently long seismograms. Today we 
first consider the eigenvalue problem in a \textbf{spherically symmetric earth model} for which the eigenfunctions can be calculated 
with relative ease. An important point is that the eigenspaces in a spherically symmetric model are \textbf{degenerate}, each having 
dimension $2l+1$ for some integer degree $l\ge0$. This material depends on the properties of \textbf{vector spherical harmonics}, but 
you are not expected to know the details. We then show using \textbf{first-order perturbation theory} how in the real Earth this degeneracy 
is \textbf{split by rotation and lateral heterogeneity}, and how the resulting first-order theory is generalised by mode coupling. To conclude, an overview is given of radial and lateral variations of density 
within the Earth, with some current research questions highlighted.

\section*{Vector spherical harmonics}
To discuss the eigenfunctions of a spherically symmetric earth model we need to review \textbf{spherical harmonic expansions}
 of scalar and vector functions. The emphasis here is only on establishing notations that will be useful later. You will 
 \textbf{not be expected to carry out any detailed calculations with either scalar or vector spherical harmonics}.

Using spherical polar co-ordinates $(r, \theta, \varphi)$ the gradient operator is given by
\begin{equation}\label{eq:1}
\nabla=\hat{\mathbf{r}}\frac{\partial}{\partial r}+\frac{1}{r}\hat{\bm{\theta}}\frac{\partial}{\partial\theta}+\frac{1}{r\sin\theta}\hat{\bm{\varphi}}\frac{\partial}{\partial\varphi},
\end{equation}
where, for example, $\hat{\mathbf{r}}$ is a unit vector in the radial direction. This expression can be conveniently written
\begin{equation}\label{eq:2}
\nabla=\hat{\mathbf{r}}\frac{\partial}{\partial r}+\frac{1}{r}\nabla_{1},
\end{equation}
which acts to define the \textbf{tangential gradient operator}, $\nabla_{1}$, for functions defined on the unit sphere, $\mathbb{S}^{2}$. 
In terms of $\nabla_{1}$ we can define related operators $\hat{\mathbf{r}}\times\nabla_{1}$ and $\nabla_{1}^{2}$. Up to suitable scale factors, the first
 is the quantum mechanical \textbf{angular momentum operator}, $\hat{\mathbf{L}}$, and the second the \textbf{total squared angular momentum}, $\hat{L}^{2}$.

Spherical harmonic functions are simultaneous eigenfunctions for the 
 \textbf{commuting} and \textbf{self-adjoint operators} $\nabla_{1}^{2}$ and $\frac{1}{\ii}\frac{\partial}{\partial\varphi}$, the latter being proportional to the \textbf{$z$-component}
  of $\hat{\mathbf{r}}\times\nabla_{1}$. The defining relations are
\begin{equation}\label{eq:3}
\nabla_{1}^{2}Y_{lm}=-l(l+1)Y_{lm}, \quad \frac{1}{\ii}\frac{\partial Y_{lm}}{\partial\varphi}=mY_{lm},
\end{equation}
where $l$ is the \textbf{degree} of the spherical harmonic which takes non-negative integer values, while for fixed $l$ the \textbf{order}, $m$, 
is an integer such that $-l\le m\le l$. Note that the eigenspace of $\nabla_{1}^{2}$ with degree $l$ is $(2l+1)$-fold degenerate, 
this being characteristic of the eigenspaces of \textbf{rotationally invariant operators}. The spherical harmonics satisfy the orthonormality relation
\begin{equation}\label{eq:4}
\int_{\mathbb{S}^{2}}Y_{lm}^{*}Y_{l'm'}\dd S=\delta_{ll'}\delta_{mm'}.
\end{equation}
They also form a \textbf{complete basis} for suitably regular functions $f$ on $\mathbb{S}^{2}$, and hence we
 can form the spherical harmonic expansion
\begin{equation}\label{eq:5}
f=\sum_{l=0}^{\infty}\sum_{m=-l}^{l}f_{lm}Y_{lm},
\end{equation}
where the expansion coefficients are
\begin{equation}\label{eq:6}
f_{lm}=\int_{\mathbb{S}^{2}}Y_{lm}^{*}f\dd S.
\end{equation}
Usually we will write such spherical harmonic expansions simply as $\sum_{lm}f_{lm}Y_{lm}$ leaving implicit the range of the degree and order. 


To work with vector fields it is convenient to introduce the \textbf{vector spherical harmonics} which are defined by
\begin{equation}\label{eq:7}
\mathbf{A}_{lm}=Y_{lm}\hat{\mathbf{r}}, \quad \mathbf{B}_{lm}=\nabla_{1}Y_{lm}, \quad \mathbf{C}_{lm}=-\hat{\mathbf{r}}\times\nabla_{1}Y_{lm}.
\end{equation}
The following orthogonality relations can be established:
\begin{align}
\int_{\mathbb{S}^{2}}\mathbf{A}_{l'm'}^{*}\cdot \mathbf{A}_{lm}\dd S&=\delta_{ll'}\delta_{mm'},\label{eq:8}\\
\int_{\mathbb{S}^{2}}\mathbf{B}_{l'm'}^{*}\cdot \mathbf{B}_{lm}\dd S&=l(l+1)\delta_{ll'}\delta_{mm'},\label{eq:9}\\
\int_{\mathbb{S}^{2}}\mathbf{C}_{l'm'}^{*}\cdot \mathbf{C}_{lm}\dd S&=l(l+1)\delta_{ll'}\delta_{mm'},\label{eq:10}
\end{align}
where $\cdot$ denotes the usual real inner product $\mathbf{a}\cdot \mathbf{b}=a_{i}b_{i}$ for vectors, while all mixed integrals, 
such as $\int_{\mathbb{S}^{2}}\mathbf{A}_{l'm'}^{*}\cdot\mathbf{C}_{lm}\dd S$, vanish. 
Note that $\mathbf{B}_{lm}$ and $\mathbf{C}_{lm}$ could be scaled such that these functions are orthonormal, but this is not necessary. 
It can be shown that the vector spherical harmonics are complete for vector fields defined on $\mathbb{S}^{2}$, this meaning 
that such a $\mathbf{u}$ can be expanded as
\begin{equation}\label{eq:11}
\mathbf{u}=\sum_{lm}(U_{lm}\mathbf{A}_{lm}+V_{lm}\mathbf{B}_{lm}+W_{lm}\mathbf{C}_{lm}),
\end{equation}
with coefficients $U_{lm}$, $V_{lm}$ and $W_{lm}$ that can be expressed in terms of $\mathbf{u}$ using the above orthogonality relations.

\subsection*{Spheroidal and toroidal modes}
In a non-rotating earth model, we recall that the eigenvalue problem takes the form
\begin{equation}\label{eq:12}
-\omega^{2}\langle \mathbf{w}|P|\mathbf{s}\rangle+\langle \mathbf{w}|H|\mathbf{s}\rangle=0,
\end{equation}
where $\omega$ is the eigenfrequency, $\mathbf{s}$ the eigenfunction, and $\mathbf{w}$ is an arbitrary test function.
When the material parameters $\rho$ and $A_{ijkl}$ are spherically symmetric,\footnote{The most general elastic tensor that is 
consistent with spherical symmetry is \textbf{transversely isotropic}
with a radial symmetry axis. There are then \textbf{five elastic moduli}, which can be taken to be Love's 
$A$, $C$, $F$, $L$ and $N$ of Lecture~14, each of which can vary with radius. In an isotropic model these reduce to the 
bulk and shear moduli $\kappa$ and $\mu$, and for simplicity we assume isotropy in what follows.} it can be shown that the eigenfunctions are of two possible types known as 
\textbf{toroidal} and \textbf{spheroidal modes}. This can be established either through a direct but lengthy calculation using 
spherical harmonics, or using representation theory for the rotation group. For reasons of both time and difficulty, we 
here just summarise the eventual results.

\begin{figure}[ht]
\centering
\includegraphics[width=\textwidth]{figures/L12F1.pdf}
\caption{Eigenfrequencies of toroidal modes within the earth model PREM of Dziewonski \&
Anderson (1981), plotted against degree $l$ for frequencies up to 50\,mHz. For each degree we see
a set of different overtones that can be labelled using the overtone number $n$, the lowest curve
being the fundamental branch $n=0$. Note also that PREM comprises a solid inner core, fluid outer
core, and solid mantle and crust. The modes seen here are mantle toroidal modes, but one could
also calculate toroidal modes confined to the inner core.}
\label{fig:1}
\end{figure}

First there are toroidal modes which have eigenfunctions of the form
\begin{equation}\label{eq:13}
\mathbf{s}(r,\theta,\varphi)=W(r)\mathbf{C}_{lm}(\theta,\varphi),
\end{equation}
with $W$ a function of radius that depends on $l$ but not $m$. 
The corresponding eigenfrequency is also independent of $m$, and hence such modes lie in $(2l+1)$-fold degenerate eigenspaces. 
It can be shown that toroidal modes are associated with no perturbation to the gravitational potential, and also that their 
eigenfunctions are non-zero only in solid parts of the earth model.

The second type of eigenfunction is the spheroidal modes that take the form
\begin{equation}\label{eq:14}
\mathbf{s}(r,\theta,\varphi)=U(r)\mathbf{A}_{lm}(\theta,\varphi)+V(r)\mathbf{B}_{lm}(\theta,\varphi),
\end{equation}
where $U$ and $V$ are functions of radius that depend on $l$ but not $m$.
Again, such an eigenfunction lies in a $(2l+1)$-fold degenerate eigenspace. Spheroidal modes 
are associated with gravitational perturbations, and their eigenfunctions can be non-zero in 
both solid and fluid regions of the earth model.

Let us consider toroidal modes further. If we put eq.~(\ref{eq:13}) into the kinetic energy form, then due to the orthogonality 
of vector spherical harmonics, we can only get a non-zero contribution from the part of the test function proportional 
to $\mathbf{C}_{lm}$. The same turns out to be true for the potential energy form, and hence it is sufficient 
to consider test functions
\begin{equation}\label{eq:15}
\mathbf{w}(r,\theta,\varphi)=w(r)\mathbf{C}_{lm}(\theta,\varphi).
\end{equation}
Using eq.~(\ref{eq:10}) it is then easy to obtain
\begin{equation}\label{eq:16}
\langle \mathbf{w}|P|\mathbf{s}\rangle=\zeta^{2}\int_{0}^{b}\rho w^{*}Wr^{2}\dd r,
\end{equation}
where $b$ is the radius of the earth model, and we have set $\zeta=\sqrt{l(l+1)}$ and assumed the earth model is entirely solid. 
In a similar manner, it can be shown that
\begin{equation}\label{eq:17}
\langle \mathbf{w}|H|\mathbf{s}\rangle=\zeta^{2}\int_{0}^{b}\mu\left(r\frac{\ddns w^{*}}{\ddns r}-w^{*}\right)\left(r\frac{\ddns W}{\ddns r}-W\right)\dd r+\zeta^{2}(\zeta^{2}-2)\int_{0}^{b}\mu w^{*}W\dd r.
\end{equation}
Note that derivation of this result is not simple, as you need to know how to differentiate vector spherical harmonics. 
Nonetheless, the end result is that the eigenvalue problem for a toroidal mode of degree $l$ reduces to the form
\begin{align}
&-\omega^{2}\int_{0}^{b}\rho w^{*}Wr^{2}\dd r+\int_{0}^{b}\mu\left(r\frac{\ddns w^{*}}{\ddns r}-w^{*}\right)\left(r\frac{\ddns W}{\ddns r}-W\right)\dd r\nonumber \\ &+(\zeta^{2}-2)\int_{0}^{b}\mu w^{*}W\dd r=0,\label{eq:18}
\end{align}
where $W$ is the radial eigenfunction, and $w$ the radial part of the test function. Through a relatively simple calculation, 
it can be shown that this is an ODE eigenvalue problem of \textbf{Sturm--Liouville type} that many of you will have seen previously.
In the case of a homogeneous planet closed-form solutions can be obtained using \textbf{spherical Bessel functions}, though such 
a solution is of little practical use. More generally, the resulting ODE eigenvalue problem can be solved numerically, or the weak 
form directly converted into an algebraic eigenvalue problem using finite element methods. For our purposes, the key result is that 
for each degree $l\ne0$ there exists a countable infinity of toroidal modes, with these different solutions being labelled by an 
\textbf{overtone number}, $n=0,1,2,\dots$, with $n=0$ being the \textbf{fundamental mode} and eigenfrequency then increasing with $n$.

For spheroidal modes a similar approach can be applied, with the eigenvalue problem reducing for each degree, $l$, 
to one involving the radial eigenfunctions $U$ and $V$. Solution of this latter problem is considerably more involved 
than that for toroidal modes, in part due to the need to model the associated gravitational potential perturbations. Nonetheless,
numerical solutions of this problem have been routinely determined for decades.

\begin{figure}[ht]
\centering
\includegraphics[width=\textwidth]{figures/L12F2.pdf}
\caption{As for Fig.~\ref{fig:1}, but here showing spheroidal mode eigenfrequencies within PREM,
including the radial modes with $l=0$. The branches are less regular than in the toroidal case, with
modes of different character (for example those whose energy is concentrated in the mantle, in the
fluid outer core, or in the inner core) crossing and interacting with one another.}
\label{fig:2}
\end{figure}



Fig.~\ref{fig:6} shows the radial eigenfunctions of a few representative modes. Several general features are worth noting. 
For toroidal modes the overtone number counts the nodes of the radial eigenfunction, so that the fundamental modes have none; 
for spheroidal modes there is no such simple rule.\footnote{A generalisation of the node-counting property to spheroidal modes does exist, formulated in terms of the so-called minor vectors of the radial equations (Woodhouse 1988; Al-Attar \& Woodhouse 2008), but it lies far beyond this course.} As the degree increases 
at fixed overtone number the eigenfunctions become confined ever closer to the surface; the fundamental modes of high degree 
are the standing-wave equivalents of surface waves, a topic that is not covered within this course. Toroidal modes of the mantle have vanishing shear traction at the core--mantle 
boundary, so that $r\,\dd W/\dd r-W$ vanishes there, while spheroidal modes penetrate the fluid outer core and the inner core, and hence 
carry information about the deep Earth. The spheroidal eigenfunctions have two radial components, $U$ for the radial displacement and 
$V$ for the tangential one; for the radial modes with $l=0$ only $U$ is present. Not all modes are so regular, and 
Fig.~\ref{fig:8} shows some of the more curious kinds: a mode dominated by a rigid translation of the inner core within the fluid 
outer core, known as the \textbf{Slichter mode}, whose period of several hours is set by gravity rather than elasticity; 
\textbf{Stoneley modes}, whose energy is trapped along a fluid--solid interface, here either the core--mantle boundary or the 
inner-core boundary, and which are the standing-wave equivalents of waves guided along such an interface; and a mode whose 
energy is concentrated within the inner core. Much of the art of free oscillation seismology lies in identifying which modes 
are sensitive to which parts of the Earth.

\begin{figure}[htbp]
\centering
\includegraphics[width=\textwidth]{figures/L12F6.pdf}
\caption{Radial eigenfunctions of some toroidal (upper row) and spheroidal (lower row) modes of PREM, each normalised to 
unit maximum amplitude. For the spheroidal modes both the radial component $U$ and the tangential component $\zeta V$, 
where $\zeta=\sqrt{l(l+1)}$, are shown. The dashed lines mark the core--mantle and inner-core boundaries. 
Calculated using the spectral-element normal-mode code of A.~Myhill.}
\label{fig:6}
\end{figure}

\begin{figure}[htbp]
\centering
\includegraphics[width=\textwidth]{figures/L12F8.pdf}
\caption{Radial eigenfunctions of some less typical spheroidal modes of PREM: the Slichter mode ${}_{1}S_{1}$, in which the inner 
core translates within the fluid outer core; Stoneley modes trapped on the core--mantle boundary and on the inner-core 
boundary; and a mode confined almost entirely to the inner core.}
\label{fig:8}
\end{figure}

We now summarise the results of this section. Within a spherically symmetric earth model there are two families of normal 
modes: toroidal and spheroidal. We can introduce an index, $q$, to distinguish these with, say, $q=1$ meaning spheroidal, 
and $q=2$ toroidal.\footnote{Along with further values if we need to distinguish toroidals in different solid parts of the 
model such as the mantle and inner core.} Having fixed the mode type, $q$, we can then consider different degrees, $l$. 
For each such degree there is an infinite set of eigenfrequencies and radial eigenfunctions labelled by their overtone number, $n$.
Finally, we then have the order, $m$, of a mode, this labelling an orthogonal basis for the $(2l+1)$-fold degenerate eigenspace. 

In this manner, we see that \textbf{four integers} $(q,l,n,m)$ are sufficient to label all the eigenfunctions of a spherically symmetric earth 
model, while the corresponding eigenfrequencies depend only on $(q,l,n)$. As a useful shorthand, we can collapse the triplet $(q,l,n)$
into a single index, $k$, and hence write $\omega_{k}$ for the eigenfrequencies and $|km\rangle$ for the eigenfunctions. 
Indeed, this is precisely the notation used in the last lecture. When needed, however, we can also write $|qlnm\rangle$ to 
explicitly characterise an eigenfunction. As a final bit of terminology, for each $(q,l,n)$ we call the eigenspace a \textbf{multiplet},
while individual elements of its basis, as labelled by $m$, are known as \textbf{singlets}. We might then, for example, talk about a 
degree two toroidal multiplet which comprises five singlets. A common labelling scheme for toroidal multiplets is to 
write ${}_{n}T_{l}$ for the multiplet of degree $l$ and overtone number $n$. Similarly, spheroidal multiplets are written ${}_{n}S_{l}$.

\subsection*{Splitting of a multiplet due to rotation and lateral heterogeneity}

Within a spherically symmetric earth model, the normal modes come in $(2l+1)$-fold degenerate multiplets, with $l$ the degree which is 
a non-negative integer. Such degeneracy is a direct consequence of the spherical symmetry of the problem, and analogous results are seen in other parts of physics. 
In particular, you know that the energy eigenspaces for the hydrogen atom are similarly degenerate, though note that there the degeneracy is larger still, owing to 
the ``hidden symmetry'' of the Coulomb problem associated with the Runge--Lenz vector. 

If such a spherically symmetric system undergoes an aspherical perturbation we expect that this degeneracy will be broken, with the singlets that
 comprise the multiplet splitting apart. This is precisely what happens in the \textbf{Zeeman effect for the hydrogen atom}, and we will apply very 
 similar mathematics to examine what rotation and lateral heterogeneities do to the free oscillations of the Earth.

We consider a spherically symmetric earth model along with the non-rotating eigenvalue problem
\begin{equation}\label{eq:19}
-\omega^{2}\langle \mathbf{w}|P_{0}|\mathbf{s}\rangle+\langle \mathbf{w}|H_{0}|\mathbf{s}\rangle=0,
\end{equation}
where we note that the subscript $0$ has been added to the kinetic and potential energy forms. A multiplet of 
this earth model is identified by the triplet $(q,l,n)$, and we write $\omega_{k}$ and $|km\rangle$ for its 
eigenfrequency and eigenfunctions using the simplified notations discussed above. We then have, in particular, that
\begin{equation}\label{eq:20}
-\omega_{k}^{2}\langle \mathbf{w}|P_{0}|km\rangle+\langle \mathbf{w}|H_{0}|km\rangle=0,
\end{equation}
for all test functions $\mathbf{w}$. 

Suppose now that we add a laterally heterogeneous perturbation to the model such that the kinetic energy form can be written
\begin{equation}\label{eq:21}
\langle \mathbf{w}|P|\mathbf{s}\rangle=\langle \mathbf{w}|P_{0}|\mathbf{s}\rangle+\epsilon\langle \mathbf{w}|P_{1}|\mathbf{s}\rangle+\cdots,
\end{equation}
where $\epsilon$ is a perturbation parameter. Using a similar notation for the potential energy form, we wish to consider the eigenvalue problem
\begin{equation}\label{eq:22}
-\omega^{2}[\langle \mathbf{w}|P_{0}|\mathbf{s}\rangle+\epsilon\langle \mathbf{w}|P_{1}|\mathbf{s}\rangle+\cdots]
+[\langle \mathbf{w}|H_{0}|\mathbf{s}\rangle+\epsilon\langle \mathbf{w}|H_{1}|\mathbf{s}\rangle+\cdots]=0.
\end{equation}
We assume that the perturbed eigenfrequency can be written
\begin{equation}\label{eq:23}
\omega=\omega_{k}+\epsilon\,\delta\omega+\cdots,
\end{equation}
while for the eigenfunctions we assume
\begin{equation}\label{eq:24}
\mathbf{s}=\sum_{m=-l}^{l}a_{m}|km\rangle+\epsilon\,\delta\mathbf{s}+\cdots,
\end{equation}
where the $a_{m}$ are coefficients to be determined. Note here that we are allowing for the zeroth-order eigenspace to be degenerate,
 this method having been applied earlier in the course when considering shear wave splitting. 
 
 Putting these expansions into the eigenvalue problem, we find at order one that
\begin{align}
&\sum_{m=-l}^{l}[-2\omega_{k}\delta\omega\langle \mathbf{w}|P_{0}|km\rangle-\omega_{k}^{2}\langle 
    \mathbf{w}|P_{1}|km\rangle+\langle \mathbf{w}|H_{1}|km\rangle]a_{m} \nonumber\\ &-\omega_{k}^{2}\langle 
        \mathbf{w}|P_{0}|\delta\mathbf{s}\rangle+\langle \mathbf{w}|H_{0}|\delta\mathbf{s}\rangle=0.\label{eq:25}
\end{align}
If we take the test function, $\mathbf{w}$, to equal $|km'\rangle$ then the terms involving $\delta\mathbf{s}$ cancel due to the self-adjointness 
 and the zeroth-order eigenvalue problem, and hence we are left with
\begin{equation}\label{eq:26}
\frac{1}{2\omega_{k}}\sum_{m=-l}^{l}[-\omega_{k}^{2}\langle km'|P_{1}|km\rangle+\langle km'|H_{1}|km\rangle]a_{m}=\delta\omega a_{m'},
\end{equation}
where we have used the orthonormality of the singlets. Such an equation is to hold for all $-l\le m'\le l$ and hence we have 
arrived at a $(2l+1)$-dimensional Hermitian eigenvalue problem for $\delta\omega$ along with the coefficients, $a_{m}$, of the 
zeroth-order eigenfunction. Assuming that the matrix elements $\langle km'|P_{1}|km\rangle$ and $\langle km'|H_{1}|km\rangle$ can 
be calculated, such an eigenvalue problem can be readily solved either in closed form for small $l$, or otherwise on a computer. 
The important qualitative conclusion of this analysis is that there will generically be $2l+1$ distinct eigenfrequency perturbations 
$\delta\omega$, this showing that the \textbf{multiplet is split into the corresponding number of non-degenerate eigenspaces}.

The above argument can be readily extended to allow for rotation of the earth model as a first-order perturbation. Indeed, we 
 generalise eq.~(\ref{eq:22}) to 
\begin{equation}\label{eq:27}
-\omega^{2}[\langle \mathbf{w}|P_{0}|\mathbf{s}\rangle+\epsilon\langle \mathbf{w}|P_{1}|\mathbf{s}\rangle+\cdots]
+\ii\omega\epsilon\langle \mathbf{w}|W|\mathbf{s}\rangle+[\langle \mathbf{w}|H_{0}|\mathbf{s}\rangle+\epsilon\langle \mathbf{w}|H_{1}|\mathbf{s}\rangle+\cdots]=0.
\end{equation}
A similar expansion to first order then leads to the Hermitian eigenvalue problem
\begin{equation}\label{eq:28}
\frac{1}{2\omega_{k}}\sum_{m=-l}^{l}[-\omega_{k}^{2}\langle km'|P_{1}|km\rangle+\ii\omega_{k}\langle km'|W|km\rangle+\langle km'|H_{1}|km\rangle]a_{m}=\delta\omega a_{m'},
\end{equation}
where we see on the left-hand side an additional set of matrix elements associated with the Coriolis force. Note 
also that terms associated with the centrifugal force need to be incorporated into the perturbation to the potential energy form.

\subsection*{The diagonal sum rule and spherically symmetric density variations}
Within the first-order splitting calculations earlier, we considered the Hermitian eigenvalue problem, eq.~(\ref{eq:28}). 
The perturbations to the singlet eigenfrequencies are equal to the eigenvalues of the matrix on the left-hand side. 
Based on so-called \textbf{selection rules} for the matrix elements,\footnote{These are nearly identical to those arising in the quantum theory
of angular momentum.} it can be shown that this matrix has zero trace. Because the trace of a matrix is invariant under
similarity transforms, it follows that the singlet eigenfrequency perturbations sum to zero. 
This result is known as the \textbf{diagonal sum rule} and was obtained by Gilbert (1971) and played a key role in early studies of
Earth structure using observations of free oscillations. 

In the previous lecture we saw that eigenfrequencies can be measured
from the amplitude spectra of sufficiently long seismograms. The first such observations were made following the 1960 Great Chilean earthquake,\footnote{``Great'' is a 
technical term used for especially large earthquakes; it does not suggest that the earthquake was a good thing.} which remains the largest event that has occurred during the instrumental period. By this time there existed good models 
of the Earth's spherically symmetric velocity structure obtained from travel time observations. There also existed plausible models 
of radial density variations that were built using knowledge of the Earth's mass, moments of inertia, and simple thermodynamic models of 
density variations in a homogeneous solid compressed under its own weight. As a consequence, the first spectra observed could be 
immediately compared to eigenfrequencies calculated in the different earth models available.

An example of this can be seen in Fig.~\ref{fig:3}. Here two reasonably clear resonance peaks can be identified, with their centre frequencies measured. 
Also shown on the plot are the predicted eigenfrequencies from two different spherical earth models for the same multiplet. It is notable, 
firstly, that these predictions are in reasonable agreement with an observed eigenfrequency, but also that while in a spherical earth 
model we would expect only one eigenfrequency in this range, there seem to be at least two in the data. The latter fact is due to 
splitting associated with rotation and lateral heterogeneity. This fact was immediately recognised, and motivated George Backus 
and Freeman Gilbert to publish the first theoretical treatment (Backus \& Gilbert 1961) of the splitting of free oscillations due to rotation. 
Indeed, as we have seen, such calculations mirror almost exactly those known from quantum mechanics for decades at that point.


\begin{figure}[ht]
\centering
\includegraphics[width=\textwidth]{figures/L12F3.png}
\caption{A part of the first free oscillation spectrum observed following the 1960 Chile
earthquake. Taken from Benioff et al.\ (1961).}
\label{fig:3}
\end{figure}

Fig.~\ref{fig:9} shows the modern equivalent, using twenty-day records of the 2011 Tohoku earthquake at three stations. Here ${}_{0}S_{2}$ 
is resolved into four of its five singlets, the central one being weak at all three stations, and ${}_{0}S_{3}$ into most of its 
seven, with the relative amplitudes of the singlets depending on the position of the station and the orientation of the source. The singlets are spread almost symmetrically 
about the degenerate frequency of PREM, and their nearly uniform spacing is the signature of the Earth's rotation, which 
splits a multiplet in much the same manner as a magnetic field splits atomic energy levels in the Zeeman effect.

\begin{figure}[htbp]
\centering
\includegraphics[width=\textwidth]{figures/L12F9.pdf}
\caption{Splitting of ${}_{0}S_{2}$ (left) and ${}_{0}S_{3}$ (right) seen in the amplitude spectra of twenty-day vertical-component 
records of the 2011 Tohoku earthquake at the stations CTAO (Australia), NNA (Peru) and BFO (Germany). The spectra are 
offset vertically, and the dashed lines mark the degenerate eigenfrequencies of the multiplets in PREM.}
\label{fig:9}
\end{figure}


During these early days of free oscillation seismology, the main emphasis was on better constraining the Earth's spherically averaged structure. 
With travel time observations this could be easily done. For many source--receiver pairs the travel time of different seismic phases 
could be measured. This information could then be reduced to plots of travel time against epicentral angle. Due to lateral heterogeneity, 
along with observational error, there was some scatter in these travel time curves, but this could be averaged out. Note that the latter 
step is not rigorously justified, but it was at least easy to do. With free oscillations, however, the situation was quite different. 
What we observe are split eigenfrequencies of the singlets, but what we really want is the eigenfrequency of the associated multiplet 
prior to this splitting. Supposing that all the singlet frequencies for a split multiplet have been measured, Gilbert's diagonal sum 
rule provides an immediate answer. We need only take the arithmetic average of the split eigenfrequencies to obtain the degenerate 
eigenfrequency corresponding to the Earth's spherically averaged structure. For some low frequency and low degree multiplets it is 
indeed possible to identify all the singlet eigenfrequencies from a single spectrum. In other cases the singlet frequencies cannot 
be seen clearly in any one spectrum, but by combining observations from many different earthquakes and stations the singlet 
frequencies can be recovered using various signal processing methods.

How do the degenerate eigenfrequencies constrain the Earth's structure? The answer is through \textbf{sensitivity kernels} 
of exactly the kind met in Lecture~20. Suppose that the density and elastic moduli of a spherically symmetric earth model are 
perturbed by small amounts $\delta\rho$, $\delta\kappa$ and $\delta\mu$. Because the eigenfrequencies are stationary with respect to 
variations of the eigenfunctions, an argument identical to that used for the first-order splitting above (with the perturbation 
now spherically symmetric, so that the multiplet is not split) shows that the change in a degenerate eigenfrequency is
\begin{equation}\label{eq:29a}
\delta\omega_{k}=\int_{0}^{b}\left(K_{\kappa}\,\delta\kappa+K_{\mu}\,\delta\mu+K_{\rho}\,\delta\rho\right)\dd r,
\end{equation}
where the kernels are quadratic in the unperturbed eigenfunctions of the mode.\footnote{This is Rayleigh's principle. Taking 
the test function equal to the eigenfunction in the perturbed eigenvalue problem, and using the unperturbed problem to 
eliminate the change in the eigenfunction, gives 
$2\omega_{k}\delta\omega_{k}=\langle k|H_{1}|k\rangle-\omega_{k}^{2}\langle k|P_{1}|k\rangle$ 
with $|k\rangle$ normalised.} For toroidal modes the kernels can be read off from eq.~(\ref{eq:18}), and with the normalisation 
$\int_{0}^{b}\rho W^{2}r^{2}\dd r=1$ they are
\begin{equation}\label{eq:29b}
K_{\mu}=\frac{1}{2\omega_{k}}\left[\left(r\frac{\ddns W}{\ddns r}-W\right)^{2}+(\zeta^{2}-2)W^{2}\right], \quad 
K_{\rho}=-\frac{\omega_{k}}{2}W^{2}r^{2}.
\end{equation}
The spheroidal kernels are more complicated, involving both $U$ and $V$ together with the perturbation to the gravitational 
potential, and we do not write them out (see Dahlen \& Tromp 1998). Fig.~\ref{fig:7} shows kernels for a few modes. Several features are important. 
The density kernel can be of either sign: added mass lowers the frequency through the kinetic energy but also increases 
the gravitational forces, and for ${}_{0}S_{2}$ the two effects compete, the kernel changing sign within the outer core. The kernels for the 
moduli are positive, as they must be, stiffer material raising the frequency. And a given mode is sensitive only where its 
eigenfunction is large, so that the high-degree fundamental modes constrain only the upper mantle, while modes such as 
${}_{0}S_{2}$ and the radial modes sample the whole planet. Because the kernels for $\kappa$, $\mu$ and $\rho$ have 
different shapes, the frequencies of a sufficiently large and varied set of modes can, in principle, be inverted for all three parameters, and 
this is how the radial density structure of the Earth is constrained by seismology.

\begin{figure}[htbp]
\centering
\includegraphics[width=\textwidth]{figures/L12F7.pdf}
\caption{Sensitivity kernels of the toroidal (upper row) and spheroidal (lower row) modes of Fig.~\ref{fig:6}. The 
kernels are shown in fractional form, so that $\delta\omega_{k}/\omega_{k}=\int_{0}^{b}(\tilde{K}_{\kappa}\,\delta\kappa/\kappa+
\tilde{K}_{\mu}\,\delta\mu/\mu+\tilde{K}_{\rho}\,\delta\rho/\rho)\dd r$. The toroidal kernels follow from eq.~(\ref{eq:29b}); 
the spheroidal ones, which include the effects of self-gravitation, were obtained numerically. The density kernel of 
${}_{0}S_{2}$ changes sign in the outer core, gravity and inertia competing for this gravest mode.}
\label{fig:7}
\end{figure}

Combining such observations with the diagonal sum rule, a large database of eigenfrequencies corresponding to the Earth's 
spherically averaged structure had been determined by the late 1970s. Taken along with travel time information, this allowed 
existing earth models to be substantially refined, and, in particular, for radial density variations to be found. The resulting 
density structure is shown in Fig.~\ref{fig:4}. These observed densities can be compared with predictions made by mineral physicists subject 
to assumptions about the Earth's interior composition. For example, it has been shown that the density of the outer core is 
significantly lower than expected if it were composed of an iron--nickel alloy as is suggested by observations of meteorites. 
This discrepancy is telling us something about how planets form and differentiate. As yet the identity of this ``light element''
in the core is uncertain, with oxygen or silicon being perhaps the most likely candidates. But what is worth emphasising is that 
such discussions are only possible due to the highly accurate constraints on the interior structure provided by seismology.

\begin{figure}[ht]
\centering
\includegraphics[width=\textwidth]{figures/L12F4.png}
\caption{Variation of density with radius in the model PREM of Dziewonski \& Anderson
(1981) as determined from the Earth's mass, moment of inertia, and over 1000
measured eigenfrequencies.}
\label{fig:4}
\end{figure}

\subsection*{Lateral density variations and mode coupling}

As noted in Lectures~19 and~20, recent work in seismology is focused on lateral heterogeneity in the mantle, this being done to 
better understand mantle convection. In Fig.~\ref{fig:5}, images from a range of S-wave velocity models are shown at different depths within 
the Earth, and a broad agreement can be seen. In particular, all models show at the base of the mantle two large regions with low 
shear wave velocity, one below the Pacific, and the other below Africa. 

\begin{figure}[ht]
\centering
\includegraphics[width=\textwidth]{figures/L12F5.png}
\caption{Depth slices from a range of different S-wave tomographic models, taken from Ritsema et al.\ (2011). 
In all cases the Pacific and African LLSVPs can be seen at the base of the mantle. The dynamical
significance of these features remains unclear.}
\label{fig:5}
\end{figure}

On the assumption that these velocity variations reflect 
temperature variations, we conclude that these so-called \textbf{Large Low Shear Velocity Provinces (LLSVPs)} are hotter than the 
surrounding mantle. If this is true, then we would also expect that these regions are lighter than the ambient mantle, and hence 
will experience an upwards buoyancy force. When LLSVPs were first identified in the early 1990s such an interpretation was taken 
for granted, and they were known as ``super plumes''. A range of subsequent work, however, has brought into question this conclusion. 
For example, joint S- and P-wave tomographic models have shown that the LLSVPs seem to be associated with a positive perturbation to 
the bulk sound speed
\begin{equation}\label{eq:30}
v_{\phi}=\sqrt{\frac{\kappa}{\rho}}.
\end{equation}
Note that this is not a speed of any elastic wave, but rather a parameter that is useful because it depends on the bulk modulus but 
not the shear modulus. 

If the lateral variations are caused by temperature, then we would expect that the shear and bulk sound velocities 
would \textbf{positively correlate}. To see this, we assume that temperature variations are sufficiently small to allow us to write
\begin{equation}\label{eq:31}
\delta\beta=\frac{\partial\beta}{\partial T}\delta T, \quad \delta v_{\phi} = \frac{\partial v_{\phi}}{\partial T}\delta T,
\end{equation}
where $\delta T$ is the temperature perturbation. The relevant partial derivatives can be estimated from mineral physics, and are 
negative for all crystalline materials, and hence
\begin{equation}\label{eq:32}
\frac{\delta v_{\phi}}{\delta\beta}=\frac{\partial v_{\phi}/\partial T}{\partial\beta/\partial T}>0,
\end{equation}
as required. The fact that shear and bulk sound speeds seem to anti-correlate in the lower mantle suggests that temperature is not 
the only cause for these lateral variations. The other possible explanation is that \textbf{compositional variations} in the mantle play an 
important role in its dynamics. To see how this might reconcile the observations, we write $X$ for a compositional variable. The 
velocity perturbations can then be expressed in terms of temperature and composition in the form
\begin{equation}\label{eq:33}
\delta\beta=\frac{\partial\beta}{\partial T}\delta T+\frac{\partial\beta}{\partial X}\delta X, \quad \delta v_{\phi}=
\frac{\partial v_{\phi}}{\partial T}\delta T+\frac{\partial v_{\phi}}{\partial X}\delta X,
\end{equation}
from which it is clear that the ratio $\delta v_{\phi}/\delta\beta$ need not always be positive. If it is accepted that 
significant compositional variations exist in the lower mantle, then we cannot determine from the tomographic studies discussed 
whether LLSVPs are light or dense, and hence the role they play in mantle dynamics. To address this latter point what 
is needed are direct estimates of their density, and it is free oscillations that have the potential to provide this information.

The first study of lateral density variations in the lower mantle using free oscillation data was undertaken in the late 1990s by 
Miaki Ishii and Jeroen Tromp, then both at Harvard (Ishii being a PhD student at the time; Ishii \& Tromp 1999). This study was based on a slight elaboration of the first-order splitting theory 
discussed earlier in the lecture, and the resulting models showed that LLSVPs are denser than the surrounding mantle. This is consistent 
with the presence of compositional variations, and hence leads to interesting questions about what these variations are, and how they were 
produced. Do, for example, LLSVPs represent a ``slab graveyard'' composed of compositionally distinct material accumulated at the base of 
the mantle through subduction? Or are they instead ancient iron-enriched features produced as the Earth differentiated into a mantle 
and core? Such questions continue to engage a wide range of Earth scientists.

But the density of LLSVPs is actually not a settled issue. 
This is because later work has shown that the simple first-order degenerate perturbation theory applied is not sufficiently accurate to 
model the effect of lateral density variations. Implicit within this theory is the assumption that on the addition of laterally varying structure
distinct multiplets do not interact with each other. This is not true, with the resulting errors being especially significant 
for nearly degenerate multiplets. A better approach is to use \textbf{normal mode coupling theory}. The idea is simple: the eigenfunctions of 
the spherical reference model form a complete basis, so any displacement field in the aspherical model, and in particular any of its 
eigenfunctions, can be expanded in this basis in the form
\begin{equation}\label{eq:34}
\mathbf{s}=\sum_{k'm'}\langle k'm'|P_{0}|\mathbf{s}\rangle|k'm'\rangle,
\end{equation}
where $|km\rangle$ are eigenfunctions in a spherical reference model. Putting such an expansion into the rotating eigenvalue problem we arrive at
\begin{equation}\label{eq:35}
\sum_{k'm'}[-\omega^{2}\langle km|P|k'm'\rangle+\ii\omega\langle km|W|k'm'\rangle+\langle km|H|k'm'\rangle]\langle k'm'|P_{0}|\mathbf{s}\rangle=0,
\end{equation}
for all $(k,m)$. This is a non-linear eigenvalue problem for the expansion coefficients $\langle k'm'|P_{0}|\mathbf{s}\rangle$ relative 
to the chosen basis. Such an equation cannot be solved directly because the matrices and vectors involved are infinite-dimensional. But it is 
possible to truncate the basis so as to obtain a finite-dimensional problem that can be solved numerically, and by taking sufficiently many 
terms in the expansion the results can be made as accurate as desired.

In practice the eigenvalue problem itself is rarely solved. What is observed is a spectrum, and this can be computed directly. 
Consider the equations of motion in the frequency domain for a given source, whose force term $\tilde{\mathbf{f}}(\omega)$ will in general 
depend on frequency, with the displacement expanded in the spherical basis as in eq.~(\ref{eq:34}), 
$\tilde{\mathbf{u}}(\omega)=\sum_{k'm'}\langle k'm'|P_{0}|\tilde{\mathbf{u}}(\omega)\rangle|k'm'\rangle$. Taking the test function equal to each basis element in turn gives
\begin{equation}\label{eq:36}
\sum_{k'm'}\left[-\omega^{2}\langle km|P|k'm'\rangle+\ii\omega\langle km|W|k'm'\rangle+\langle km|H|k'm'\rangle\right]
\langle k'm'|P_{0}|\tilde{\mathbf{u}}(\omega)\rangle=\langle km|\tilde{\mathbf{f}}(\omega)\rangle,
\end{equation}
which for each frequency is a linear system for the expansion coefficients, the same matrices appearing as in eq.~(\ref{eq:35}). 
Splitting the kinetic and potential energy forms into their spherically symmetric parts and the perturbations, and recalling that the 
basis is orthonormal with respect to $P_{0}$ and diagonalises $H_{0}$, this system takes the form
\begin{equation}\label{eq:37}
(\omega_{k}^{2}-\omega^{2})\langle km|P_{0}|\tilde{\mathbf{u}}(\omega)\rangle+\sum_{k'm'}A_{km,k'm'}(\omega)\langle k'm'|P_{0}|\tilde{\mathbf{u}}(\omega)\rangle
=\langle km|\tilde{\mathbf{f}}(\omega)\rangle,
\end{equation}
where the matrix $A$ collects the matrix elements of $P_{1}$, $W$ and $H_{1}$, all of which are small. For a given frequency 
$\omega$, the diagonal terms $\omega_{k}^{2}-\omega^{2}$ are large except for the modes whose eigenfrequencies lie close to 
$\omega$, and so the system is dominated by its diagonal apart from a small block associated with those modes. This structure 
means that the system can be solved efficiently by modern iterative methods, at a cost of a few multiplications by $A$ 
per frequency, without ever forming the inverse or diagonalising the full coupling matrix. This is the 
\textbf{direct solution method} of Al-Attar et al.\ (2012), and it is now the standard means of computing 
long-period spectra in laterally heterogeneous earth models. In practice the system is solved not at real frequencies but along 
a line shifted slightly into the complex plane. This keeps the diagonal terms away from zero, and corresponds, once the solutions are 
transformed back to the time domain, to multiplying the seismogram by a decaying exponential, which is harmless over a record of finite length.

The first-order splitting theory is recovered by keeping only the $2l+1$ singlets of a single multiplet in the expansion, 
and this so-called \textbf{self-coupling} approximation is the workhorse of most splitting studies. Within it the effect of lateral heterogeneity 
on a multiplet is summarised by its \textbf{splitting function}, a function on the sphere whose spherical harmonic coefficients 
are, up to known factors, the matrix elements in eq.~(\ref{eq:28}). It can be thought of as a map of the 
local perturbation to the degenerate eigenfrequency, obtained by integrating the lateral heterogeneity against the mode's 
radial sensitivity kernels. Splitting functions are measured from the spectra of many earthquakes, and then form the data 
for inversions for lateral structure, including density. Their appeal is that they collapse a great many spectra into a few 
coefficients per multiplet that depend only on Earth structure. Their weakness is that they are defined 
within the self-coupling approximation, and their sensitivity to density is weak and easily contaminated by errors in that 
approximation, particularly for multiplets that lie close in frequency to others. Only in the past decade or so has it 
become possible to compute the spectra of realistic 3D earth models in the full-coupling manner just described, 
and such calculations have shown that self-coupling is not accurate enough for density (Akbarashrafi et al.\ 2018), the 
errors it introduces being comparable to the density signal itself. The problem is more severe still for modes 
whose energy is concentrated near the core--mantle boundary, exactly those most sensitive to the density of the LLSVPs. 
Hence the results of Ishii and Tromp cannot be taken at face value, and work continues on making full-coupling calculations 
efficient enough to use routinely within inversions. This being said, a more recent but related study was
led by Harriet Lau using observations of solid Earth tides (Lau et al.\ 2017). In this case sophisticated and accurate numerical modelling 
was done, with the result of this analysis still supporting the idea that LLSVPs are dense stable features 
at the base of the mantle. 

Finally, it is worth drawing together the threads of the last few lectures. Once spectra can be computed accurately for a 3D earth 
model, the natural next step is to fit them directly, in the manner of the waveform inversions of Lecture~20, rather than 
proceeding through intermediate measurements such as splitting functions. The ingredient needed is the sensitivity of the 
spectra to the model, and this is provided by the adjoint method: applied to the direct solution method, 
it yields the kernels of the misfit with respect to Earth structure at the cost of one further solution per frequency 
(Adourian et al.\ 2024). Full-waveform inversion of the Earth's free oscillations for its long-wavelength 
structure, including density, is thus now feasible, and this is where much of the current work in the subject lies.

\section*{What you need to know and be able to do}
\begin{enumerate}
    \item[(i)] How eigenfunctions of a spherically symmetric model can be classified in terms of mode type, degree, overtone number, and order. 
    You do not need to learn the definition of vector spherical harmonics, and will not be asked to perform detailed calculations with them.
    \item[(ii)] How to solve the toroidal mode eigenvalue problem in a homogeneous earth model using spherical Bessel functions (see the 
    worked examples for this lecture, available on the TiS). Note that you might be asked to perform such calculations in the exam, but would be provided with the necessary 
    information about special functions.
    \item[(iii)] How degenerate perturbation theory can be used to understand mode splitting due to rotation and lateral heterogeneity.     
    \item[(iv)] What the diagonal sum rule says (but not the derivation), and how it was applied in determining spherically averaged 
    density variations.
    \item[(v)] The idea that convection in the mantle might be driven in part by compositional variations, and how this complicates 
    the dynamic interpretation of tomographic velocity models.
    \item[(vi)] The idea of full coupling, namely that the eigenfunctions or the forced response of a laterally heterogeneous earth 
    model can be expanded in the modes of a spherical reference model, and why the simpler first-order splitting theory can be inadequate, 
    particularly for density. You do not need the details of how such calculations are done, nor anything about splitting 
    functions or adjoint methods for spectra.
\end{enumerate}
